% Abschließende Folgerungen aus den früh eingeführten Strukturaxiomen.
\chapter{Eindeutigkeit und Wiederverwendung der Strukturen}

Die Wortaxiome wurden unmittelbar nach dem Aufbau der elementaren
Anfügung eingeführt, die Baumaxiome unmittelbar nach dem Nachweis der
Konstruktor- und Trägereigenschaften. Die Wortinduktion steht durch W3
axiomatisch bereit; die Bauminduktion und die Rekursionssätze sind aus
den jeweiligen Strukturaxiomen bewiesen. Abschließend
werden ihre Konsequenzen für den Darstellungswechsel und die
Wiederverwendung zusammengeführt.

\section{Eindeutigkeit der Strukturen und das leere Alphabet}

Ein \emph{strukturtreuer Isomorphismus} ist hier eine Bijektion, welche
den Anfang und die Anfügung beziehungsweise sämtliche Blattbeschriftungen
und das geordnete Zusammenfügen erhält. Das Alphabet wird dabei
festgehalten; seine Buchstaben werden nicht umbenannt.

\FormulaThmDeltaK[Eindeutiger Isomorphismus freier Wortstrukturen]{%
  \begin{aligned}[t]
    &\mathsf{WortStr}_A(W,e,s)\dsep
      \mathsf{WortStr}_A(V,d,r)\vdash{}\\[-2pt]
    &\exists!j\;\bigl(j\colon W\bij V\land
      \mathsf{WRec}_A(j;W,e,s;V,d,r)\bigr)
  \end{aligned}
}{WordStructureUniqueIsomorphism}{%
  \DeltaRow{Mengen}{A\dsep W\dsep V\dsep e\dsep d}
  \DeltaRow{Funktionen}{s\dsep r\dsep h\dsep h'\dsep j}}
\begin{tabproofwide}
  \proofstepwidestar[1]{\mathsf{WortStr}_A(W,e,s)}{%
    \rA}
  \proofstepwidestar[2]{\mathsf{WortStr}_A(V,d,r)}{%
    \rA}
  \proofstepwidestar[2]{d\in V\land r\colon V\times A\to V}{%
    \FormulaRefAuto{WordStructureTypingAxiom}[ax]{2}}
  \proofstepwidestar[1,2]{\exists!j\;\mathsf{WRec}_A(j;W,e,s;V,d,r)}{%
    \FormulaRefAuto{WordStructureRecursion}[thm]{1,\rAEa{3},\rAEb{3}}}
  \proofstepwidestar[1,2]{\exists j\;\mathsf{WRec}_A(j;W,e,s;V,d,r)}{%
    \FormulaRefAuto{\exists! x\,P(x)\vdash\exists x\,P(x)}[thm]{4}}
  \proofstepwidestar[6]{\mathsf{WRec}_A(j;W,e,s;V,d,r)}{%
    \rA}
  \proofstepwidestar[1,2,6]{j\colon W\bij V}{%
    \FormulaRefAuto{WordStructureRecursorIsomorphism}[thm]{1,2,6}}
  \proofstepwidestar[1,2,6]{\exists j\;(j\colon W\bij V\land\mathsf{WRec}_A(j;W,e,s;V,d,r))}{%
    \rEI{\rAI{7,6}}}
  \proofstepwidestar[1,2]{\exists j\;(j\colon W\bij V\land\mathsf{WRec}_A(j;W,e,s;V,d,r))}{%
    \rEE{5,6,8}}
  \proofstepwidestar[10]{j\colon W\bij V\land\mathsf{WRec}_A(j;W,e,s;V,d,r)}{%
    \rA}
  \proofstepwidestar[11]{h\colon W\bij V\land\mathsf{WRec}_A(h;W,e,s;V,d,r)}{%
    \rA}
  \proofstepwidestar[1,2,10,11]{j=h}{%
    \FormulaRefAuto{WordRecursionTwoSolutionsEqual}[thm]{1,\rAEa{3},\rAEb{3},\rAEb{10},\rAEb{11}}}
  \proofstepwidestar[1,2]{\exists!j\;(j\colon W\bij V\land\mathsf{WRec}_A(j;W,e,s;V,d,r))}{%
    \UEI{9,10,11,12}}
\end{tabproofwide}

\FormulaThmDeltaK[Eindeutiger Isomorphismus freier Baumstrukturen]{%
  \begin{aligned}[t]
    &\mathsf{BaumStr}_A(T,\ell,k)\dsep
      \mathsf{BaumStr}_A(V,m,n)\vdash{}\\[-2pt]
    &\exists!j\;\bigl(j\colon T\bij V\land
      \mathsf{TRec}_A(j;T,\ell,k;V,m,n)\bigr)
  \end{aligned}
}{BinaryTreeStructureUniqueIsomorphism}{%
  \DeltaRow{Mengen}{A\dsep T\dsep V}
  \DeltaRow{Funktionen}{\ell\dsep k\dsep m\dsep n\dsep h\dsep h'\dsep j}}
\begin{tabproofwide}
  \proofstepwidestar[1]{\mathsf{BaumStr}_A(T,\ell,k)}{%
    \rA}
  \proofstepwidestar[2]{\mathsf{BaumStr}_A(V,m,n)}{%
    \rA}
  \proofstepwidestar[2]{m\colon A\to V\land n\colon V\times V\to V}{%
    \FormulaRefAuto{BinaryTreeStructureTypingAxiom}[ax]{2}}
  \proofstepwidestar[1,2]{\exists!j\;\mathsf{TRec}_A(j;T,\ell,k;V,m,n)}{%
    \FormulaRefAuto{BinaryTreeStructureRecursion}[thm]{1,\rAEa{3},\rAEb{3}}}
  \proofstepwidestar[1,2]{\exists j\;\mathsf{TRec}_A(j;T,\ell,k;V,m,n)}{%
    \FormulaRefAuto{\exists! x\,P(x)\vdash\exists x\,P(x)}[thm]{4}}
  \proofstepwidestar[6]{\mathsf{TRec}_A(j;T,\ell,k;V,m,n)}{%
    \rA}
  \proofstepwidestar[1,2,6]{j\colon T\bij V}{%
    \FormulaRefAuto{BinaryTreeStructureRecursorIsomorphism}[thm]{1,2,6}}
  \proofstepwidestar[1,2,6]{\exists j\;(j\colon T\bij V\land\mathsf{TRec}_A(j;T,\ell,k;V,m,n))}{%
    \rEI{\rAI{7,6}}}
  \proofstepwidestar[1,2]{\exists j\;(j\colon T\bij V\land\mathsf{TRec}_A(j;T,\ell,k;V,m,n))}{%
    \rEE{5,6,8}}
  \proofstepwidestar[10]{j\colon T\bij V\land\mathsf{TRec}_A(j;T,\ell,k;V,m,n)}{%
    \rA}
  \proofstepwidestar[11]{h\colon T\bij V\land\mathsf{TRec}_A(h;T,\ell,k;V,m,n)}{%
    \rA}
  \proofstepwidestar[1,10,11]{j=h}{%
    \FormulaRefAuto{BinaryTreeRecursionTwoSolutionsEqual}[thm]{1,\rAEb{10},\rAEb{11}}}
  \proofstepwidestar[1,2]{\exists!j\;(j\colon T\bij V\land\mathsf{TRec}_A(j;T,\ell,k;V,m,n))}{%
    \UEI{9,10,11,12}}
\end{tabproofwide}

\Needspace{12\baselineskip}
\FormulaThmDeltaK[Freie Strukturen über dem leeren Alphabet]{%
  \begin{aligned}[t]
    &\text{(i)}\quad\mathsf{WortStr}_{\varnothing}(W,e,s)\vdash W=\{e\},\\[-2pt]
    &\text{(ii)}\quad\mathsf{BaumStr}_{\varnothing}(T,\ell,k)\vdash T=\varnothing.
  \end{aligned}
}{FreeStructuresEmptyAlphabet}{%
  \DeltaRow{Mengen}{W\dsep T\dsep e\dsep u\dsep a\dsep x\dsep y}\DeltaRow{Funktionen}{s\dsep\ell\dsep k}}
\begin{tabproofsplitwide}
 \proofpartwideindR[Wortstruktur]{\mathsf{WortStr}_{\varnothing}(W,e,s)\vdash W=\{e\}}
 {\mathsf{WortStr}_{\varnothing}(W,e,s)\vdash W=\{e\}}[WordStructureEmptyAlphabet]
  \proofstepwidestar[1]{\mathsf{WortStr}_{\varnothing}(W,e,s)}{%
    \rA}
  \proofstepwidestar[1]{e\in W}{%
    \rAEa{\FormulaRefAuto{WordStructureTypingAxiom}[ax]{1}}}
  \proofstepwidestar[1]{\{e\}\subseteq W}{%
    \FormulaRefAuto{x\in A\vdash\{x\}\subseteq A}[thm]{2}}
  \proofstepwidestar[]{e\in\{e\}}{%
    \FormulaRefAuto{x\in\{x\}}[thm]}
  \proofstepwidestar[5]{u\in\{e\}}{%
    \rA}
  \proofstepwidestar[6]{a\in\varnothing}{%
    \rA}
  \proofstepwidestar[]{a\notin\varnothing}{%
    \FormulaRefAuto{x\not\in\varnothing}[thm]}
  \proofstepwidestar[6]{s(u,a)\in\{e\}}{%
    \FormulaRefAuto{A\dsep\neg A\vdash B}[thm]{6,7}}
  \proofstepwidestar[]{\forall u\in\{e\}\,\forall a\in\varnothing\;s(u,a)\in\{e\}}{%
    \rUI{\rRI{5,\rUI{\rRI{6,8}}}}}
  \proofstepwidestar[1]{\{e\}=W}{%
    \FormulaRefAuto{WordStructureMinimalityAxiom}[ax]{1,3,4,9}}
  \proofstepwidestar[1]{W=\{e\}}{%
    \FormulaRefAuto{a=b\vdash b=a}[thm]{10}}
 \closeproofpartwideindR
 \Needspace{14\baselineskip}
 \proofpartwideindR[Baumstruktur]{\mathsf{BaumStr}_{\varnothing}(T,\ell,k)\vdash T=\varnothing}
 {\mathsf{BaumStr}_{\varnothing}(T,\ell,k)\vdash T=\varnothing}[BinaryTreeStructureEmptyAlphabet]
  \proofstepwidestar[1]{\mathsf{BaumStr}_{\varnothing}(T,\ell,k)}{%
    \rA}
  \proofstepwidestar[1]{\ell\colon\varnothing\to T}{%
    \rAEa{\FormulaRefAuto{BinaryTreeStructureTypingAxiom}[ax]{1}}}
  \proofstepwidestar[]{\varnothing\subseteq T}{%
    \FormulaRefAuto{\varnothing\subseteq A}[thm]}
  \proofstepwidestar[1]{\ell[\varnothing]=\varnothing}{%
    \FormulaRefAuto{F\colon A\to B\vdash F[\varnothing]=\varnothing}[thm]{2}}
  \proofstepwidestar[]{\varnothing\subseteq\varnothing}{%
    \FormulaRefAuto{A\subseteq A}[thm]}
  \proofstepwidestar[1]{\ell[\varnothing]\subseteq\varnothing}{%
    \rIE{\FormulaRefAuto{a=b\vdash b=a}[thm]{4},5}}
  \proofstepwidestar[7]{x\in\varnothing}{%
    \rA}
  \proofstepwidestar[8]{y\in\varnothing}{%
    \rA}
  \proofstepwidestar[]{x\notin\varnothing}{%
    \FormulaRefAuto{x\not\in\varnothing}[thm]}
  \proofstepwidestar[7]{k(x,y)\in\varnothing}{%
    \FormulaRefAuto{A\dsep\neg A\vdash B}[thm]{7,9}}
  \proofstepwidestar[]{\forall x,y\in\varnothing\;k(x,y)\in\varnothing}{%
    \rUI{\rRI{7,\rUI{\rRI{8,10}}}}}
  \proofstepwidestar[1]{\varnothing=T}{%
    \FormulaRefAuto{BinaryTreeStructureMinimalityAxiom}[ax]{1,3,6,11}}
  \proofstepwidestar[1]{T=\varnothing}{%
    \FormulaRefAuto{a=b\vdash b=a}[thm]{12}}
 \closeproofpartwideindR
\end{tabproofsplitwide}

Der Unterschied ist beabsichtigt: Ein Wort darf leer sein; ein endlicher
voller Klammerungsbaum besitzt ein Blatt. Ohne Blattbeschriftungen gibt es
daher keinen solchen Baum. Die Rekursionssätze gelten auch in diesen
Grenzfällen. Bei Wörtern bleibt der vorgegebene Anfangswert erforderlich;
bei der leeren Baumstruktur ist die gesuchte Funktion die leere Funktion.

\section{Wiederverwendung in Algebra und formaler Syntax}

Die Axiome fassen den Band zu einer wiederverwendbaren Grundlage zusammen.
Die axiomatische Wortinduktion liefert die Existenz der Zerlegung;
Trennung und Injektivität der Konstruktoren sichern die Eindeutigkeit.
Bei Bäumen liefert die Minimalität das entsprechende Induktionsprinzip.
Die Rekursionssätze wurden direkt
aus den Strukturaxiomen bewiesen; auf ihnen beruht der eindeutige Transport. Ein späterer Beweis darf
deshalb die Strukturgesetze benutzen, ohne die Positionen eines Wortes
oder den Wortcode eines Baumes erneut auszuwerten.

Beispielsweise erhält die Baumrekursion mit
\(f(a)=\LetterWord a\) und \(g(u,v)=\WordConcat u v\) das Blattwort in
jeder freien Baumstruktur. Mit einer beliebigen binären Operation auf
\(X\) erhält man ihre Auswertung. Die Eindeutigkeit der Rekursion zeigt,
dass diese Funktionen unter dem eindeutigen Isomorphismus mit den früher
definierten Blattwort- und Auswertungsfunktionen übereinstimmen. Die
bereits bewiesenen Aussagen zu Klammerungen bleiben damit in jeder
Darstellung gültig.

Für Band~48, \emph{Axiomatische Mengenlehre II}, ist dieselbe Trennung
zwischen Konstruktion und Strukturgesetzen nützlich. Eine Formel besitzt
einen äußeren Konstruktor und dessen unmittelbare Teilformeln. Die
entsprechenden Regeln für Syntax müssen allerdings die tatsächlichen
Stelligkeiten berücksichtigen: Negation hat eine Teilformel, eine binäre
Verknüpfung zwei, und ein Quantor zusätzlich einen Variablenparameter.
Atomare Formeln tragen ihre eigenen Parameter. Die hier behandelte
Binärbaumstruktur beschreibt unmittelbar die volle binäre
Klammerungsstruktur; sie ist noch kein vollständiges Modell dieser
verschiedenen Formelkonstruktoren.

Für eine solche Syntax wird das vorliegende Muster auf die betreffende
Familie von Konstruktoren übertragen: Jeder Konstruktor ist auf seinen
Parametern und Teilobjekten injektiv, die Bilder verschiedener
Konstruktoren sind disjunkt, und jede unter allen Konstruktoren
abgeschlossene Teilmenge ist die gesamte Syntaxmenge. Nach Konstruktion
eines passenden Modells sind Strukturinduktion und Strukturrekursion
mit je einem Fall pro Konstruktor zu beweisen. Dabei dürfen Identifikationen
wie Gleichheit bis zur Umbenennung gebundener Variablen erst in einem
gesonderten Schritt eingeführt werden; sie verändern die Aussage der
eindeutigen Zerlegung.

Die Verbindung mit den Graphen- und Baumbänden liefert dazu die
geometrische Sicht: Vorkommen von Teilobjekten sind Knoten, die Beziehung
zum unmittelbaren Teilobjekt gibt gerichtete Kanten, und das Vergessen
ihrer Richtung gibt den zugrunde liegenden ungerichteten Baum.
Verschiedene Vorkommen gleicher Teilobjekte bleiben dabei verschiedene
Knoten. Die hier bewiesenen Strukturgesetze erklären, welche Information
die geordnete und beschriftete Baumdarstellung für spätere rekursive
Definitionen bewahren muss.
